% GENERATED FILE. Edit canonical lessons, then run npm run generate:books.
% canonical-source-hash: fnv1a64:823777e6e7ffa19a
% canonical-lessons: ML-C320-prasastamaya, ML-C320-sadharana, ML-C320-priyappetta, ML-C320-bhayankaramaya, ML-C320-atipoli

\chapter{Famous, Ordinary, Dear, Terrible, Awesome}
\label{ch:ml-famous-ordinary-dear-terrible-awesome}
\hlchaptermodality{320}

\begin{chapteropening}
\textbf{By the end of this chapter:} \emph{I can say famous, ordinary, dear, terrible and awesome.}

Complete the last lesson of chapter 320: I can say famous, ordinary, dear, terrible and awesome.
\end{chapteropening}

\section[\texorpdfstring{\ml{പ്രശസ്തമായ} (praśastamāya)}{പ്രശസ്തമായ (praśastamāya)}]{\ml{പ്രശസ്തമായ} (praśastamāya) — famous}

\label{lesson:ML-C320-prasastamaya}



\begin{quote}
Before the new one: say the Malayalam for arrogance, then the Malayalam for strict.
\end{quote}

\subsection*{You'll want to know: \ml{പ്രശസ്തമായ}}
\textbf{\ml{പ്രശസ്തമായ}} — \emph{praśastamāya} — \textquotedblleft{}famous\textquotedblright{}.

\begin{cousinweb}
\textbf{\ml{പ്രശസ്തമായ}} (\emph{praśastamāya}) joins Sanskrit \emph{praśasta}, \textquotedblleft{}praised, celebrated\textquotedblright{}, to native \emph{āya}, \textquotedblleft{}having become\textquotedblright{} (DEDR 333). DEDR 333 compares Sanskrit \emph{ām}.
\end{cousinweb}

\begin{grammarlens}[title={What you've built}]
One more word for this chapter.
\end{grammarlens}

\subsection*{Your turn}
\begin{itemize}
\raggedright
  \item \emph{Say it:} \emph{praśastamāya}
  \item \emph{Say it:} \emph{praśastamāya}, once more
  \item \emph{Say it:} say \emph{praśastamāya}
  \item \emph{Recall:} say the Malayalam for arrogance, then the Malayalam for strict, then say \emph{praśastamāya} again
\end{itemize}

\begin{tcolorbox}[breakable,colback=teal!4,colframe=teal!35!black,title={Before you move on}]
What does \ml{പ്രശസ്തമായ} mean? (\textquotedblleft{}Famous\textquotedblright{}.) Say it once more.
\end{tcolorbox}
\section[\texorpdfstring{\ml{സാധാരണ} (sādhāraṇa)}{സാധാരണ (sādhāraṇa)}]{\ml{സാധാരണ} (sādhāraṇa) — ordinary, usual}

\label{lesson:ML-C320-sadharana}



\begin{quote}
Before the new one: say the Malayalam for strict, then the Malayalam for famous.
\end{quote}

\subsection*{You'll want to know: \ml{സാധാരണ}}
\textbf{\ml{സാധാരണ}} — \emph{sādhāraṇa} — \textquotedblleft{}ordinary, usual\textquotedblright{}.

\begin{cousinweb}
\textbf{\ml{സാധാരണ}} (\emph{sādhāraṇa}) $\leftarrow$ Sanskrit \emph{sādhāraṇa}, \textquotedblleft{}common, general\textquotedblright{}, taken over whole as a tatsama.
\end{cousinweb}

\begin{grammarlens}[title={What you've built}]
One more word for this chapter.
\end{grammarlens}

\subsection*{Your turn}
\begin{itemize}
\raggedright
  \item \emph{Say it:} \emph{sādhāraṇa}
  \item \emph{Say it:} \emph{sādhāraṇa}, once more
  \item \emph{Say it:} say \emph{sādhāraṇa}
  \item \emph{Recall:} say the Malayalam for strict, then the Malayalam for famous, then say \emph{sādhāraṇa} again
\end{itemize}

\begin{tcolorbox}[breakable,colback=teal!4,colframe=teal!35!black,title={Before you move on}]
What does \ml{സാധാരണ} mean? (\textquotedblleft{}Ordinary, usual\textquotedblright{}.) Say it once more.
\end{tcolorbox}
\section[\texorpdfstring{\ml{പ്രിയപ്പെട്ട} (priyappeṭṭa)}{പ്രിയപ്പെട്ട (priyappeṭṭa)}]{\ml{പ്രിയപ്പെട്ട} (priyappeṭṭa) — dear, favourite}

\label{lesson:ML-C320-priyappetta}



\begin{quote}
Before the new one: say the Malayalam for famous, then the Malayalam for ordinary.
\end{quote}

\subsection*{You'll want to know: \ml{പ്രിയപ്പെട്ട}}
\textbf{\ml{പ്രിയപ്പെട്ട}} — \emph{priyappeṭṭa} — \textquotedblleft{}dear, favourite\textquotedblright{}.

\begin{cousinweb}
\textbf{\ml{പ്രിയപ്പെട്ട}} (\emph{priyappeṭṭa}) joins Sanskrit \emph{priya}, \textquotedblleft{}beloved\textquotedblright{} (CDIAL 8974), to native \emph{peṭṭa}, a form of the verb \textquotedblleft{}to fall into, become\textquotedblright{} (DEDR 3853).
\end{cousinweb}

\begin{grammarlens}[title={What you've built}]
One more word for this chapter.
\end{grammarlens}

\subsection*{Your turn}
\begin{itemize}
\raggedright
  \item \emph{Say it:} \emph{priyappeṭṭa}
  \item \emph{Say it:} \emph{priyappeṭṭa}, once more
  \item \emph{Say it:} say \emph{priyappeṭṭa}
  \item \emph{Recall:} say the Malayalam for famous, then the Malayalam for ordinary, then say \emph{priyappeṭṭa} again
\end{itemize}

\begin{tcolorbox}[breakable,colback=teal!4,colframe=teal!35!black,title={Before you move on}]
What does \ml{പ്രിയപ്പെട്ട} mean? (\textquotedblleft{}Dear, favourite\textquotedblright{}.) Say it once more.
\end{tcolorbox}
\section[\texorpdfstring{\ml{ഭയങ്കരമായ} (bhayaṅkaramāya)}{ഭയങ്കരമായ (bhayaṅkaramāya)}]{\ml{ഭയങ്കരമായ} (bhayaṅkaramāya) — terrible, awful}

\label{lesson:ML-C320-bhayankaramaya}



\begin{quote}
Before the new one: say the Malayalam for ordinary, then the Malayalam for dear.
\end{quote}

\subsection*{You'll want to know: \ml{ഭയങ്കരമായ}}
\textbf{\ml{ഭയങ്കരമായ}} — \emph{bhayaṅkaramāya} — \textquotedblleft{}terrible, awful\textquotedblright{}.

\begin{cousinweb}
\textbf{\ml{ഭയങ്കരമായ}} (\emph{bhayaṅkaramāya}) is Sanskrit \emph{bhayaṅkara}, \textquotedblleft{}causing fear\textquotedblright{}: \emph{bhaya}, \textquotedblleft{}fear\textquotedblright{}, and \emph{kara}, \textquotedblleft{}making\textquotedblright{}, with native \emph{āya}, \textquotedblleft{}having become\textquotedblright{} (DEDR 333). DEDR 333 compares Sanskrit \emph{ām}.
\end{cousinweb}

\begin{grammarlens}[title={What you've built}]
One more word for this chapter.
\end{grammarlens}

\subsection*{Your turn}
\begin{itemize}
\raggedright
  \item \emph{Say it:} \emph{bhayaṅkaramāya}
  \item \emph{Say it:} \emph{bhayaṅkaramāya}, once more
  \item \emph{Say it:} say \emph{bhayaṅkaramāya}
  \item \emph{Recall:} say the Malayalam for ordinary, then the Malayalam for dear, then say \emph{bhayaṅkaramāya} again
\end{itemize}

\begin{tcolorbox}[breakable,colback=teal!4,colframe=teal!35!black,title={Before you move on}]
What does \ml{ഭയങ്കരമായ} mean? (\textquotedblleft{}Terrible, awful\textquotedblright{}.) Say it once more.
\end{tcolorbox}
\section[\texorpdfstring{\ml{അടിപൊളി} (aṭipoḷi)}{അടിപൊളി (aṭipoḷi)}]{\ml{അടിപൊളി} (aṭipoḷi) — awesome, great (slang)}

\label{lesson:ML-C320-atipoli}



\begin{quote}
Before the new one: say the Malayalam for dear, then the Malayalam for terrible.
\end{quote}

\subsection*{You'll want to know: \ml{അടിപൊളി}}
\textbf{\ml{അടിപൊളി}} — \emph{aṭipoḷi} — \textquotedblleft{}awesome, great (slang)\textquotedblright{}.

\begin{cousinweb}
\textbf{\ml{അടിപൊളി}} (\emph{aṭipoḷi}) is slang built from native parts: \emph{aṭi}, \textquotedblleft{}a blow\textquotedblright{} (DEDR 77), and \emph{poḷi}, from the verb \textquotedblleft{}to break open\textquotedblright{} (DEDR 4560), much like English \emph{smashing}.
\end{cousinweb}

\begin{grammarlens}[title={What you've built}]
One more word for this chapter.
\end{grammarlens}

\subsection*{Your turn}
\begin{itemize}
\raggedright
  \item \emph{Say it:} \emph{aṭipoḷi}
  \item \emph{Say it:} \emph{aṭipoḷi}, once more
  \item \emph{Say it:} say \emph{aṭipoḷi}
  \item \emph{Recall:} say the Malayalam for dear, then the Malayalam for terrible, then say \emph{aṭipoḷi} again
\end{itemize}

\begin{tcolorbox}[breakable,colback=teal!4,colframe=teal!35!black,title={Before you move on}]
What does \ml{അടിപൊളി} mean? (\textquotedblleft{}Awesome, great (slang)\textquotedblright{}.) Say it once more.
\end{tcolorbox}
